\documentclass[11pt]{article} % For LaTeX2e
\usepackage{pdt_report}
\input{commands}

% It is strongly recommended to use hyperref, especially for readibility.
\definecolor{boilermakergold}{HTML}{cfb991}
\definecolor{steel}{HTML}{555960}
\definecolor{coolgray}{HTML}{6f727b}
\usepackage[pagebackref,breaklinks,colorlinks,allcolors=coolgray]{hyperref}

% Add your title here.
\title{%
  A \LaTeX{} Template for Report and Publications
}

% The \author macro works with any number of authors. There are two commands
% used to separate the names and addresses of multiple authors: \And and \AND.
%
% Using \And between authors leaves it to LaTeX to determine where to break the
% lines. Using \AND forces a line break at that point. So, if LaTeX puts 3 of 4
% authors names on the first line, and the last on the second line, try using
% \AND instead of \And before the third author name.
\author{%
  Juanwu~Lu~$^{1}$\thanks{Use footnote for providing further information
    about author (webpage, alternative address)---but \emph{not} for acknowledging funding agencies.}\\
  College of Engineering, Purdue University \\
  West Lafayette, IN 47907, USA \\
  \texttt{juanwu@purdue.edu}
  % examples of more authors
  % \And
  % Coauthor \\
  % Affiliation \\
  % Address \\
  % \texttt{email} \\
  % \AND
  % Coauthor \\
  % Affiliation \\
  % Address \\
  % \texttt{email} \\
  % \And
  % Coauthor \\
  % Affiliation \\
  % Address \\
  % \texttt{email} \\
  % \And
  % Coauthor \\
  % Affiliation \\
  % Address \\
  % \texttt{email} \\
}

% Change the title of your supplementary material here.
\renewcommand{\appendixtocname}{Supplementary Material}
\renewcommand{\appendixpagename}{Supplementary Material}

\begin{document}

\maketitle

% -----------------------------------------
% Main contents
\begin{abstract}
    \label{sec: abstract}
    \lipsum[1]
\end{abstract}
\keywords{First keyword \and Second keyword \and More}

\input{contents/1_introduction}
\section{Related Works}
\label{sec: review}

\paragraph{Paragraph Heading 1.}
\label{para: heading_1}
\lipsum[4]

\paragraph{Paragraph Heading 2.}
\label{para: heading_2}
\lipsum[5]

\section{Method}
\label{sec: method}

\subsection{Level-2 Heading}
\label{subsec: heading}
\lipsum[5]
\lipsum[6]
\begin{figure}[!ht]
    \centering
    \begin{tikzpicture}
        % Nodes
        \node[obs] (oprev) {$\vo_{t-1}$};
        \node[obs, right=of oprev] (ocurr) {$\vo_{t}$};
        \node[obs, right=of ocurr] (onext) {$\vo_{t+1}$};
        \node[latent, below=of oprev] (sprev) {$\vs_{t-1}$};
        \node[latent, right=of sprev] (scurr) {$\vs_{t}$};
        \node[const, left=of sprev] (spast) {$\ldots\quad$};
        \node[latent, right=of scurr] (snext) {$\vs_{t+1}$};
        \node[const, right=of snext] (sfutr) {$\quad\ldots$};
        \node[obs, below=of sprev, xshift=0.8cm] (aprev) {$\va_{t-1}$};
        \node[obs, below=of scurr, xshift=0.8cm] (acurr) {$\va_{t}$};

        % Edges
        \edge {spast} {sprev};
        \edge {sprev} {oprev, scurr, aprev};
        \edge {scurr} {ocurr, snext, acurr};
        \edge {snext} {onext, sfutr};
        \edge {aprev} {scurr};
        \edge {acurr} {snext};
    \end{tikzpicture}
    \caption{Example figure caption.}
    \label{fig: bayesnet}
\end{figure}


\subsubsection{Level-3 Heading}
\label{subsubsec: heading}
\lipsum[7]
\lipsum[8]
\begin{equation}
    \gJ(\vtheta)\triangleq{
        \mathbb{E}_{\tau\sim{p^{\pi}_{\vtheta}(\tau)}}
        \left[\sum\limits_{t=1}^{\left|\tau\right|}r(s_{t},a_{t})\right]
    }
\end{equation}

\section{Experiments}
\label{sec: experiment}
\lipsum[9]
\lipsum[10]

\begin{table}
  \caption{Sample table title}
  \label{sample-table}
  \centering
  \begin{tabular}{@{}lll@{}}
    \toprule
    \multicolumn{2}{c}{Part}                   \\
    \cmidrule(r){1-2}
    Name     & Description     & Size ($\mu$m) \\
    \midrule
    Dendrite & Input terminal  & $\sim$100     \\
    Axon     & Output terminal & $\sim$10      \\
    Soma     & Cell body       & up to $10^6$  \\
    \bottomrule
  \end{tabular}
\end{table}

\section{Conclusion}
\label{sec: conclusion}
\lipsum[11]
\lipsum[12]

\begin{ack}
Use unnumbered first level headings for the acknowledgments. All acknowledgments
go at the end of the paper before the list of references. Moreover, you are required to declare
funding (financial activities supporting the submitted work) and competing interests (related financial activities outside the submitted work).
\end{ack}

\section*{References}


References follow the acknowledgments in the camera-ready paper. Use unnumbered first-level heading for
the references. Any choice of citation style is acceptable as long as you are
consistent. It is permissible to reduce the font size to \verb+small+ (9 point)
when listing the references.
Note that the Reference section does not count towards the page limit.
\medskip


{%
}


% -----------------------------------------
% Appendix
\newpage
\appendix
\begin{appendices}
    % 1. Start a new local contents list
    \startcontents[appendices]

    % 2. Print the local table of contents (only for what follows)
    \printcontents[appendices]{l}{1}{\setcounter{tocdepth}{2}}
    \newpage

    \section{First Appendix Section}
    \label{appx: first}
    \lipsum[12]
    \lipsum[13]

    \section{Second Appendix Section}
    \label{appx: second}
    \lipsum[14]
    \lipsum[15]

\end{appendices}


\end{document}
